\documentclass{meetingminutes}
\meetingcommittee{สภามหาวิทยาลัยราชภัฏอุตรดิตถ์}
\meetingnumber{๔}{๒๕๖๘}
\meetingdate{วันที่ ๔ เมษายน พ.ศ. ๒๕๖๘}
\meetingplace{ห้องประชุมสภามหาวิทยาลัยราชภัฏอุตรดิตถ์}
\meetingstart{๑๓.๐๐ น.}
\meetingend{๑๖.๐๐ น.}

\begin{document}
\maketitle

\psection{ผู้มาประชุม}
\begin{participants}
  \pmember{1}{นายกสภามหาวิทยาลัยราชภัฏอุตรดิตถ์}{มหาวิทยาลัยราชภัฏอุตรดิตถ์}{ประธาน}
  \pmember{2}{อธิการบดีมหาวิทยาลัยราชภัฏอุตรดิตถ์}{มหาวิทยาลัยราชภัฏอุตรดิตถ์}{กรรมการ}
  \pmember{3}{กรรมการสภามหาวิทยาลัย (ผู้ทรงคุณวุฒิ)}{มหาวิทยาลัยราชภัฏอุตรดิตถ์}{กรรมการ}
  \pmember{4}{ผู้แทนคณะเทคโนโลยีอุตสาหกรรม}{คณะเทคโนโลยีอุตสาหกรรม}{กรรมการ}
\end{participants}

\psection{ผู้ไม่มาประชุม}
\noindent ไม่มี\par

\startmeeting
\openingchair{นายกสภามหาวิทยาลัยราชภัฏอุตรดิตถ์}{กล่าวเปิดการประชุมและดำเนินการประชุมตามระเบียบวาระ ดังนี้}

% ============================================================
\agenda{๑}{เรื่องแจ้งเพื่อทราบ}
% ============================================================

ประธานแจ้งให้ที่ประชุมทราบว่า หลักสูตรวิศวกรรมศาสตรมหาบัณฑิต สาขาวิชาวิศวกรรมคอมพิวเตอร์และปัญญาประดิษฐ์
คณะเทคโนโลยีอุตสาหกรรม ได้ผ่านกระบวนการพิจารณาจากหน่วยงานที่เกี่ยวข้องทุกระดับมาแล้ว ได้แก่
คณะกรรมการบัณฑิตศึกษา (๔ มีนาคม ๒๕๖๗) สภาวิชาการ (๒๑ มีนาคม ๒๕๖๗)
และคณะกรรมการประจำคณะเทคโนโลยีอุตสาหกรรม (๑๑ ธันวาคม ๒๕๖๗)
จึงนำเสนอสภามหาวิทยาลัยเพื่อพิจารณาอนุมัติในการประชุมครั้งนี้

\noindent\textbf{มติที่ประชุม}\quad ที่ประชุมรับทราบ

% ============================================================
\agenda{๒}{พิจารณาอนุมัติหลักสูตรวิศวกรรมศาสตรมหาบัณฑิต
สาขาวิชาวิศวกรรมคอมพิวเตอร์และปัญญาประดิษฐ์}
% ============================================================

ที่ประชุมสภามหาวิทยาลัยพิจารณาหลักสูตรวิศวกรรมศาสตรมหาบัณฑิต สาขาวิชาวิศวกรรมคอมพิวเตอร์และปัญญาประดิษฐ์
ที่ผ่านกระบวนการกลั่นกรองจากทุกระดับมาเรียบร้อยแล้ว

ที่ประชุมอภิปรายและพิจารณาสาระสำคัญของหลักสูตร ประกอบด้วย

\begin{thailist}
  \item ผลลัพธ์การเรียนรู้ที่คาดหวัง (PLO) จำนวน ๕ ข้อ ครอบคลุมสมรรถนะด้านวิศวกรรมคอมพิวเตอร์และปัญญาประดิษฐ์
  ที่ตอบสนองความต้องการของตลาดแรงงานและชุมชนในภาคเหนือตอนล่าง
  \item โครงสร้างหลักสูตร ๒ แผน (แผน ก วิทยานิพนธ์ และแผน ข สารนิพนธ์)
  มีหน่วยกิตรวมตรงตามเกณฑ์มาตรฐาน
  \item อาจารย์ประจำหลักสูตร ๗ คน มีคุณสมบัติตรงตามเกณฑ์ กมอ. พ.ศ. ๒๕๖๕
  \item แผนการรับนักศึกษาและกลยุทธ์การประชาสัมพันธ์หลักสูตรแก่ผู้มีส่วนได้ส่วนเสีย
\end{thailist}

ที่ประชุมพิจารณาแล้วเห็นว่าหลักสูตรมีความสมบูรณ์และพร้อมสำหรับการเปิดสอน

\resolution{ที่ประชุมมีมติอนุมัติหลักสูตรวิศวกรรมศาสตรมหาบัณฑิต
สาขาวิชาวิศวกรรมคอมพิวเตอร์และปัญญาประดิษฐ์ คณะเทคโนโลยีอุตสาหกรรม
มหาวิทยาลัยราชภัฏอุตรดิตถ์ ให้เปิดสอนได้ตั้งแต่ปีการศึกษา ๒๕๖๘
และให้เผยแพร่ผลลัพธ์การเรียนรู้ที่คาดหวังต่อผู้มีส่วนได้ส่วนเสียผ่านช่องทางที่เหมาะสม}

% ============================================================
\agenda{๓}{เรื่องอื่น ๆ}
% ============================================================

ไม่มีวาระเพิ่มเติม

\noindent\textbf{มติที่ประชุม}\quad ที่ประชุมรับทราบ

\adjournmeeting

\begin{sigblock}
\typist{ผู้แทนคณะเทคโนโลยีอุตสาหกรรม}
\reviewer{กรรมการสภามหาวิทยาลัย}{กรรมการ}
\chairsig{นายกสภามหาวิทยาลัยราชภัฏอุตรดิตถ์}{ประธานที่ประชุม}
\end{sigblock}

\end{document}
